\section{Fazit}
Der Einsatz von KI-gestützten Werkzeugen durch Angreifer erleichtert sowohl die Identifikation als auch die Ausnutzung von Schwachstellen erheblich.

Insbesondere in den Bereichen Phishing und Social Engineering ermöglichen Large Language Models (LLMs) eine deutlich effizientere und schnellere Durchführung von Angriffen. Dabei ist für eine hinreichende Qualität kein leistungsstarkes Spitzenmodell erforderlich. Da moderne Software zunehmend sicherer gestaltet wird und ein Großteil der erfolgreichen Angriffe ihren Ursprung in Phishing- oder Social-Engineering-Kampagnen hat, senkt der Einsatz von LLMs die Einstiegshürde für derartige Angriffsvektoren weiter ab. Dabei verändert ein LLM den grundlegenden Angriffsablauf nicht zwingend, sondern optimiert einzelne Teilschritte - etwa die zielgruppengerechte Übersetzung und sprachliche Anpassung von Inhalten. Das daraus resultierende Risiko trifft Privatpersonen in stärkerem Maße als Unternehmen, da im privaten Umfeld deutlich weniger technische und organisatorische Schutzmaßnahmen etabliert sind als im Unternehmensbereich.


Besonders relevante Angriffsvektoren sind die automatisierte Schwachstellenanalyse von Quellcode und physischen Geräten. Diese Vektoren zeichnen sich dadurch aus, dass der Angreifer den gesamten Ablauf kontrolliert und nicht auf menschliches Fehlverhalten angewiesen ist - ein wesentlicher Unterschied zu sozialen Angriffsvektoren wie Phishing oder Social Engineering. Empirische Ergebnisse, etwa aus Experimenten mit Claude Opus, belegen, dass auf diese Weise auch kritische Schwachstellen identifiziert werden können, wie etwa eine FreeBSD-Schwachstelle mit einem CVSS-Score von 10.0. Im Rahmen des eigenen Selbstversuchs konnten für physische Geräte hingegen keine vergleichbar gravierenden Ergebnisse erzielt werden. %TODO: Prüfen, ob Studien zu ähnlichen Befunden gelangen.

Grundsätzlich ist die Generierung funktionsfähiger Exploits auch mit öffentlich zugänglichen Sprachmodellen möglich. Kritisch anzumerken ist dabei, dass der Zugang zu Claude Opus auf bestimmte Nutzergruppen beschränkt ist, weshalb die entsprechenden Ergebnisse mit Vorbehalt zu interpretieren sind. Studienergebnisse zeigen, dass öffentliche Modelle in ihrer Leistungsfähigkeit nur geringfügig hinter proprietären Systemen wie Claude Opus zurückbleiben - bei deutlich niedrigeren Kosten. Sollte die von Anthropic kommunizierte Fähigkeit zutreffen, dass Claude Opus eigenständig Schwachstellen in Projekten identifizieren kann, ohne explizite Nutzerinteraktion vorauszusetzen, würde dies ein erhebliches offensives Potenzial darstellen.

Zusammenfassend lässt sich festhalten, dass KI-gestützte Werkzeuge sowohl auf Angreifer- als auch auf Verteidigungsseite zunehmend an Bedeutung gewinnen und in beiden Kontexten strategisch eingesetzt werden sollten. Es ist davon auszugehen, dass sich hieraus eine asymmetrische Wettbewerbsdynamik entwickelt - insbesondere im Bereich öffentlich zugänglicher Codebases stellt sich die Frage, welche Partei über ausreichend Ressourcen verfügt, um KI-Modelle effektiver zur Codeanalyse einzusetzen.
